For \(|S_c|=2\), let the two active coefficients be \(L_c/2\) and \(-L_c/2\). Relabeling the active arms and multiplying the published two-arm posterior-mean rule by \(L_c/2\) yields an exact finite-\(n\) saddle under independent probability-\(1/2\) assignment, which equals \(q^\star\); the lifted least-favorable prior gives the matching converse, and \[d_n(c)=C_0(c)C_A n^{-4/3}+o(n^{-4/3}).\] For \(c=c^\dagger\), every \(n,M\geq1\) admits a computable orbit procedure and rational prior satisfying \[B_n(\nu_{n,M})\leq\rho_n^\dagger\leq U_{n,M}\leq\Lambda_{n,M}^\dagger\leq B_n(\nu_{n,M})+\frac{1}{4M^2}.\] Hence, for every positive normalization \(a_n\) and integer sequence \(M_n\) with \(a_n/M_n^2\to0\), \[a_n\{U_{n,M_n}-B_n(\nu_{n,M_n})\}\longrightarrow0.\] If \(a_nd_n(c^\dagger)\) is subsequently proved to converge, these finite-\(n\) upper and lower certificate improvements converge to the same limit. This is the strongest available attainment statement: it is exact and computable at finite \(n\), but it does not identify \(a_n\), establish second-order optimality of \(q^\star\), or produce a continuum feedback rule for \(|S_c|\geq3\).